\documentclass[11pt,a4paper]{article}
\usepackage[utf8]{inputenc}
\usepackage[french]{babel}
\usepackage{amsmath,amssymb,amsthm}
\usepackage{geometry}
\geometry{hmargin=2.5cm,vmargin=3cm}
\usepackage{tikz}
\usetikzlibrary{cd,positioning,shapes}
\usepackage{listings}
\usepackage{xcolor}

\lstset{
    language=Python,
    basicstyle=\ttfamily\small,
    keywordstyle=\color{blue}\bfseries,
    stringstyle=\color{red},
    commentstyle=\color{gray}\itshape,
    numbers=left,
    numberstyle=\tiny\color{gray},
    breaklines=true,
    frame=single,
    showstringspaces=false,
    literate={é}{{\'e}}1 {è}{{\`e}}1 {à}{{\`a}}1 {ê}{{\^e}}1 {î}{{\^i}}1 {ç}{{\c{c}}}1 {ï}{{\"i}}1 {É}{{\'E}}1
}

\title{\Huge LIVRABLE IA T6 : FORMALISATION DE LA DIVERGENCE DE KULLBACK-LEIBLER \\ \large Cours magistral, exercices corrigés et travaux pratiques}
\author{\Large Charles EDOU NZE}
\date{\today}

\begin{document}
\maketitle
\tableofcontents
\newpage
\part{Cours Magistral}




\section{Formalisation Mathématique de l'Information}

La divergence de Kullback-Leibler (ou entropie relative) est une mesure asymétrique de la dissimilitude entre deux distributions de probabilités $P$ et $Q$. Issu de la théorie de l'information (Shannon, 1948), ce concept quantifie l'excès d'information ou la pénalité asymptotique subie lorsque la distribution $Q$ est utilisée pour approximer la véritable distribution $P$.

En intelligence artificielle, la divergence de Kullback-Leibler est fondamentale. Elle sert de fonction objectif dans la minimisation de l'erreur d'approximation lors de l'entraînement de modèles génératifs (tels que les auto-encodeurs variationnels ou les modèles de diffusion) et constitue le fondement mathématique de la fonction de perte 	extit{Cross-Entropy}.

\section{Formalisation}

Soit $(\mathcal{X}, \mathcal{F}, \lambda)$ un espace mesuré (généralement $\mathbb{R}^n$ avec la mesure de Lebesgue). Soient $P$ et $Q$ deux mesures de probabilité sur cet espace.

\subsection{Définition via les densités}

Supposons que $P$ et $Q$ admettent des densités $p$ et $q$ par rapport à $\lambda$.


\subsection{Exemples Concrets Immédiats}

\textbf{Exemple 1 : Divergence entre deux pièces de monnaie (Variables de Bernoulli)}
Soit $P$ la loi d'une pièce biaisée avec $p=0.8$ (pile) et $1-p=0.2$ (face). Soit $Q$ la loi d'une pièce équilibrée $q=0.5$.
$$ D_{KL}(P || Q) = 0.8 \ln\left(\frac{0.8}{0.5}\right) + 0.2 \ln\left(\frac{0.2}{0.5}\right) \approx 0.8 \times 0.470 + 0.2 \times (-0.916) = 0.376 - 0.183 = 0.193 $$
L'asymétrie est visible si on calcule $D_{KL}(Q || P)$ :
$$ D_{KL}(Q || P) = 0.5 \ln\left(\frac{0.5}{0.8}\right) + 0.5 \ln\left(\frac{0.5}{0.2}\right) = 0.5 \times (-0.470) + 0.5 \times 0.916 = 0.223 $$
Ainsi, $D_{KL}(P || Q) \neq D_{KL}(Q || P)$.

\textbf{Exemple 2 : Divergence nulle pour des distributions identiques}
Si $P = Q$ (par exemple $p = 0.5, q = 0.5$), alors :
$$ D_{KL}(P || P) = 0.5 \ln(1) + 0.5 \ln(1) = 0 $$

\textbf{Exemple 3 : Divergence infinie (support non inclus)}
Si $P(X=1) = 1$ et $Q(X=1) = 0$, alors :
$$ D_{KL}(P || Q) = 1 \ln\left(\frac{1}{0}\right) = +\infty $$
Cela illustre l'absolue continuité requise : $P$ doit être absolument continue par rapport à $Q$ ($P \ll Q$).

\textbf{Exemple 4 : Deux lois uniformes}
Soient $P = \mathcal{U}([0, 1])$ et $Q = \mathcal{U}([0, 2])$.
Pour $x \in [0, 1]$, $p(x) = 1$ et $q(x) = 0.5$.
$$ D_{KL}(P || Q) = \int_{0}^{1} 1 \ln\left(\frac{1}{0.5}\right) dx = \ln(2) \approx 0.693 $$
À l'inverse, pour $x \in ]1, 2]$, $p(x) = 0$ et $q(x) = 0.5$. Comme $p(x)>0$ sur un ensemble où $q(x)=0$ n'arrive pas, mais pour $D_{KL}(Q || P)$, on intègre sur $[0, 2]$. Sur $]1, 2]$, $p(x) = 0$, d'où $D_{KL}(Q || P) = +\infty$.

\textbf{Exemple 5 : Divergence KL entre deux Gaussiennes univariées}
Soient $P = \mathcal{N}(\mu_1, \sigma_1^2)$ et $Q = \mathcal{N}(\mu_2, \sigma_2^2)$.
La formule analytique rigoureuse donne :
$$ D_{KL}(P || Q) = \ln\left(\frac{\sigma_2}{\sigma_1}\right) + \frac{\sigma_1^2 + (\mu_1 - \mu_2)^2}{2\sigma_2^2} - \frac{1}{2} $$
Si $\mu_1 = 0, \sigma_1 = 1$ et $\mu_2 = 2, \sigma_2 = 1$ :
$$ D_{KL}(P || Q) = \ln(1) + \frac{1 + 4}{2} - 0.5 = 2.0 $$
Cette forme quadratique par rapport aux moyennes explique l'efficacité de la KL pour approcher des gaussiennes.


> \textbf{Définition (Divergence de Kullback-Leibler) :}
> On définit la divergence KL de $Q$ par rapport à $P$ par l'intégrale de Lebesgue :
> $$D_{KL}(P \| Q) = \int_{\mathcal{X}} p(x) \ln\left( \frac{p(x)}{q(x)} \right) d\lambda(x)$$
> \textit{Condition :} On suppose que $P$ est absolument continue par rapport à $Q$ ($P \ll Q$), c'est-à-dire que $q(x)=0 \implies p(x)=0$. Sinon, la divergence est $+\infty$.

\subsection{Propriétés Fondamentales}

> \textbf{Inégalité de Gibbs (Positivité) :}
> Pour toutes distributions de probabilité $P$ et $Q$ :
> $$D_{KL}(P \| Q) \ge 0$$
> Avec égalité $D_{KL}(P \| Q) = 0$ si et seulement si $P = Q$ presque partout.

> \textbf{Asymétrie :} Attention, en général $D_{KL}(P \| Q) \neq D_{KL}(Q \| P)$. Ce n'est donc pas une "distance" au sens métrique du terme (Jalon 51).

\section{Démonstrations}

\subsection{Démonstration de la positivité (via l'inégalité de Jensen)}

1. \textbf{Réécriture :} $D_{KL}(P \| Q) = \int p(x) [-\ln(q(x)/p(x))] d\lambda(x) = \mathbb{E}_P [ -\ln(Q/P) ]$.
2. \textbf{Utilisation de la convexité :} La fonction $\phi(u) = -\ln(u)$ est strictement convexe sur $]0, +\infty[$.
3. \textbf{Inégalité de Jensen :} Pour toute variable aléatoire $U$ et fonction convexe $\phi$ :
   $$\mathbb{E}[\phi(U)] \ge \phi(\mathbb{E}[U])$$
4. \textbf{Application :} Posons $U = q(X)/p(X)$ où $X \sim P$.
   $$D_{KL}(P \| Q) = \mathbb{E}_P [ \phi(U) ] \ge \phi(\mathbb{E}_P[U])$$
5. \textbf{Calcul de l'espérance de U :}
   $$\mathbb{E}_P[U] = \int p(x) \frac{q(x)}{p(x)} d\lambda(x) = \int q(x) d\lambda(x) = 1$$ (car $Q$ est une probabilité).
6. \textbf{Conclusion :}
   $D_{KL}(P \| Q) \ge \phi(1) = -\ln(1) = 0$.
   La positivité est démontrée.

\section{Exercices d'Application}

\subsection{Exercice 1 : KL entre deux Gaussiennes}
\textbf{Énoncé :} Calculer $D_{KL}(P \| Q)$ pour $P = \mathcal{N}(\mu_1, \sigma^2)$ et $Q = \mathcal{N}(\mu_2, \sigma^2)$.
\textbf{Correction Détaillée :}
1. \textbf{Log-ratio :} $\ln(p(x)/q(x)) = \frac{1}{2\sigma^2} [ (x-\mu_2)^2 - (x-\mu_1)^2 ]$.
2. \textbf{Développement :} $(x-\mu_2)^2 - (x-\mu_1)^2 = x^2 - 2x\mu_2 + \mu_2^2 - (x^2 - 2x\mu_1 + \mu_1^2) = 2x(\mu_1 - \mu_2) + \mu_2^2 - \mu_1^2$.
3. \textbf{Intégration par rapport à P :} $\mathbb{E}_P[x] = \mu_1$.
   $D_{KL} = \frac{1}{2\sigma^2} [ 2\mu_1(\mu_1 - \mu_2) + \mu_2^2 - \mu_1^2 ] = \frac{1}{2\sigma^2} [ \mu_1^2 - 2\mu_1\mu_2 + \mu_2^2 ]$.
4. \textbf{Résultat :} $D_{KL}(P \| Q) = \frac{(\mu_1 - \mu_2)^2}{2\sigma^2}$.
La divergence augmente avec le carré de la distance entre les moyennes.

\subsection{Exercice 2 : Niveau Avancé (Lien avec l'Entropie Croisée)}
\textbf{Énoncé :} Montrer que minimiser la Cross-Entropy entre des données et un modèle revient à minimiser la divergence KL.
\textbf{Correction Détaillée :}
$H(P, Q) = -\int p(x) \ln q(x) dx$.
On remarque que $D_{KL}(P \| Q) = \int p \ln p - \int p \ln q = -H(P) + H(P, Q)$.
Comme l'entropie des données $H(P)$ est constante par rapport aux paramètres du modèle, minimiser $H(P, Q)$ est équivalent à minimiser $D_{KL}$.

\section{Application en Intelligence Artificielle}

- \textbf{Le Pont Théorique :} La divergence KL est la fonction de perte par défaut pour tous les modèles probabilistes. Elle fait le lien entre la \textbf{Théorie de la Mesure} et la \textbf{Théorie de l'Information}.
- \textbf{Example Concret :}
    - \textbf{Variational Auto-Encoders (VAE) :} La fonction de coût est la somme d'une erreur de reconstruction et d'un terme KL qui force la distribution latente à être proche d'une Gaussienne standard.
    - \textbf{Apprentissage par Renforcement (PPO) :} L'algorithme Proximal Policy Optimization utilise une contrainte KL pour éviter que la nouvelle politique ne s'éloigne trop de l'ancienne, garantissant ainsi une mise à jour stable.
    - \textbf{Classification Multi-classe :} Le calcul de la perte Softmax est rigoureusement une minimisation de la divergence KL entre la distribution "one-hot" des étiquettes et les probabilités prédites par le réseau.

\section{Liens Sémantiques}

- \textbf{Concepts Précédents requis :} Jalon 71 (Théorèmes de Fubini-Tonelli).md, Jalon 66 (Construction de l'intégrale de Lebesgue pour les fonctions mesurables positives.).md
- \textbf{Concepts Futurs dépendants :} Jalon 85 (Axiomes de Kolmogorov).md, Jalon 140 (Classifieur de Bayes optimal).md

\newpage
\part{Exercices d\'Application et de Concours}
\subsection*{Exercice 1 : Divergence KL entre deux variables de Bernoulli \quad $\bigstar\star\star$}
\subsubsection*{Énoncé}
Soient $P$ et $Q$ deux distributions de Bernoulli de paramètres $p$ et $q$.
Démontrez la positivité stricte de $D_{KL}(P||Q)$ pour $p \neq q$.
\subsubsection*{Correction}
Par définition, $D_{KL}(P||Q) = p \ln\left(\frac{p}{q}\right) + (1-p) \ln\left(\frac{1-p}{1-q}\right)$.
L'inégalité log de Gibbs $x \ln(x/y) \ge x - y$ implique que cette somme est $\ge (p-q) + ((1-p)-(1-q)) = 0$.
L'égalité n'est atteinte que si $p=q$.

\subsection*{Exercice 2 : Asymétrie de la divergence KL \quad $\bigstar\star\star$}
\subsubsection*{Énoncé}
Avec $P \sim \mathcal{B}(0.1)$ et $Q \sim \mathcal{B}(0.9)$, calculez $D_{KL}(P||Q)$ et $D_{KL}(Q||P)$.
\subsubsection*{Correction}
$D_{KL}(P||Q) = 0.1 \ln(1/9) + 0.9 \ln(9) \approx -0.22 + 1.97 = 1.75$.
$D_{KL}(Q||P) = 0.9 \ln(9) + 0.1 \ln(1/9) \approx 1.75$, ici la symétrie est accidentelle due à $p = 1-q$.
Prenons $P \sim \mathcal{B}(0.2), Q \sim \mathcal{B}(0.5)$ :
$D_{KL}(P||Q) = 0.2 \ln(0.4) + 0.8 \ln(1.6) \approx -0.18 + 0.37 = 0.19$.
$D_{KL}(Q||P) = 0.5 \ln(2.5) + 0.5 \ln(0.625) \approx 0.45 - 0.23 = 0.22$.
L'asymétrie est claire.

\subsection*{Exercice 3 : KL entre lois normales (même variance) \quad $\bigstar\bigstar\star$}
\subsubsection*{Énoncé}
Calculez $D_{KL}(\mathcal{N}(\mu_1, \sigma^2) || \mathcal{N}(\mu_2, \sigma^2))$.
\subsubsection*{Correction}
La formule donne $\ln(1) + \frac{\sigma^2 + (\mu_1-\mu_2)^2}{2\sigma^2} - 0.5 = \frac{(\mu_1-\mu_2)^2}{2\sigma^2}$.
La divergence est proportionnelle au carré de la distance euclidienne des moyennes, divisé par la variance.

\subsection*{Exercice 4 : Entropie croisée et divergence KL \quad $\bigstar\bigstar\star$}
\subsubsection*{Énoncé}
Montrez que $H(P, Q) = H(P) + D_{KL}(P||Q)$.
\subsubsection*{Correction}
$H(P, Q) = -\sum P(x) \ln Q(x) = -\sum P(x) \ln P(x) + \sum P(x) \ln\left(\frac{P(x)}{Q(x)}\right) = H(P) + D_{KL}(P||Q)$.
Minimiser l'entropie croisée revient à minimiser la KL par rapport à $Q$ si $P$ est fixe.

\subsection*{Exercice 5 : Inégalité de Pinsker (introduction) \quad $\bigstar\bigstar\bigstar$}
\subsubsection*{Énoncé}
Rappelez l'inégalité de Pinsker liant la variation totale et la divergence KL.
\subsubsection*{Correction}
L'inégalité s'écrit $\delta(P, Q) \le \sqrt{\frac{1}{2} D_{KL}(P||Q)}$. Elle contrôle l'erreur absolue de probabilité globale.

\subsection*{Exercice 6 : KL pour des lois exponentielles \quad $\bigstar\bigstar\star$}
\subsubsection*{Énoncé}
Calculez $D_{KL}(Exp(\lambda_1) || Exp(\lambda_2))$.
\subsubsection*{Correction}
$\int_0^\infty \lambda_1 e^{-\lambda_1 x} \ln\left(\frac{\lambda_1 e^{-\lambda_1 x}}{\lambda_2 e^{-\lambda_2 x}}\right) dx$.
$= \ln(\frac{\lambda_1}{\lambda_2}) + \frac{\lambda_2 - \lambda_1}{\lambda_1}$.

\subsection*{Exercice 7 : Divergence KL et vraisemblance \quad $\bigstar\bigstar\bigstar$}
\subsubsection*{Énoncé}
Montrez que la maximisation de la log-vraisemblance équivaut à minimiser la divergence KL empirique.
\subsubsection*{Correction}
Soit $P_{emp}$ la distribution empirique. $D_{KL}(P_{emp} || Q_\theta) = \sum P_{emp}(x) \ln\frac{P_{emp}(x)}{Q_\theta(x)}$.
Minimiser ceci par rapport à $\theta$ revient à maximiser $\sum P_{emp}(x) \ln Q_\theta(x) = \frac{1}{N} \sum_{i=1}^N \ln Q_\theta(x_i)$, qui est la log-vraisemblance.

\subsection*{Exercice 8 : Divergence de Jensen-Shannon \quad $\bigstar\bigstar\bigstar$}
\subsubsection*{Énoncé}
Définissez la divergence de Jensen-Shannon et vérifiez sa symétrie.
\subsubsection*{Correction}
$JSD(P||Q) = \frac{1}{2} D_{KL}(P || M) + \frac{1}{2} D_{KL}(Q || M)$ où $M = \frac{1}{2}(P+Q)$. Par construction, elle est symétrique.

\subsection*{Exercice 9 : Propriété d'additivité \quad $\bigstar\bigstar\bigstar$}
\subsubsection*{Énoncé}
Montrez que $D_{KL}(P_1 \times P_2 || Q_1 \times Q_2) = D_{KL}(P_1 || Q_1) + D_{KL}(P_2 || Q_2)$.
\subsubsection*{Correction}
Le logarithme du produit des probabilités est la somme des logarithmes. L'intégrale double se sépare par linéarité et par le fait que les marginales intègrent à 1.

\subsection*{Exercice 10 : Borne inférieure de Cramer-Rao \quad $\bigstar\bigstar\bigstar\bigstar$}
\subsubsection*{Énoncé}
En utilisant la divergence KL, démontrez que l'information de Fisher quantifie la courbure locale de l'espace des distributions.
\subsubsection*{Correction}
Par développement de Taylor de $D_{KL}(P_\theta || P_{\theta + d\theta})$, le terme d'ordre 1 est nul, et le terme d'ordre 2 correspond à la matrice d'information de Fisher.

\newpage
\part{Travaux Pratiques et Simulations Algorithmiques}
\subsection*{TP 1 : Calcul de KL Bernoulli}
\begin{lstlisting}[language=Python]
import math
def kl_bernoulli(p, q):
    if p == 0: return math.log(1 / (1 - q))
    if p == 1: return math.log(1 / q)
    return p * math.log(p / q) + (1 - p) * math.log((1 - p) / (1 - q))

assert abs(kl_bernoulli(0.5, 0.5)) < 1e-6
assert kl_bernoulli(0.8, 0.5) > 0
\end{lstlisting}

\subsection*{TP 2 : KL Gaussienne}
\begin{lstlisting}[language=Python]
import math
def kl_gaussian(mu1, sig1, mu2, sig2):
    return math.log(sig2 / sig1) + (sig1**2 + (mu1 - mu2)**2) / (2 * sig2**2) - 0.5

assert abs(kl_gaussian(0, 1, 0, 1)) < 1e-6
assert kl_gaussian(0, 1, 2, 1) == 2.0
\end{lstlisting}

\subsection*{TP 3 : Asymétrie}
\begin{lstlisting}[language=Python]
import math
def kl_bernoulli(p, q):
    return p * math.log(p / q) + (1 - p) * math.log((1 - p) / (1 - q))

pq = kl_bernoulli(0.2, 0.5)
qp = kl_bernoulli(0.5, 0.2)
assert abs(pq - qp) > 1e-6
\end{lstlisting}

\subsection*{TP 4 : Entropie croisée}
\begin{lstlisting}[language=Python]
import math
def cross_entropy(p, q):
    return -(p * math.log(q) + (1 - p) * math.log(1 - q))

def entropy(p):
    return -(p * math.log(p) + (1 - p) * math.log(1 - p))

def kl(p, q):
    return p * math.log(p / q) + (1 - p) * math.log((1 - p) / (1 - q))

ce = cross_entropy(0.3, 0.6)
h_plus_kl = entropy(0.3) + kl(0.3, 0.6)
assert abs(ce - h_plus_kl) < 1e-6
\end{lstlisting}

\subsection*{TP 5 : Jensen-Shannon}
\begin{lstlisting}[language=Python]
import math
def kl(p, q):
    return p * math.log(p / q) + (1 - p) * math.log((1 - p) / (1 - q))

def jsd(p, q):
    m = 0.5 * (p + q)
    return 0.5 * kl(p, m) + 0.5 * kl(q, m)

assert abs(jsd(0.2, 0.8) - jsd(0.8, 0.2)) < 1e-6
\end{lstlisting}

\end{document}